\chapter{Das Kreuz als Geleite}
\quellseite{10}

Als Friedrich Wilhelm Weber, ein im Jahre 1894 verstorbener namhafter Dichter und Arzt in Westphalen, vom Vaterhause Abschied nahm, um als Student nach Greifswald zu ziehen, schnitzte er sich im nahen Walde ein kleines Kreuz. Dies kleine Kreuz war sein Begleiter durchs Leben, und die erstarrte Hand des 80jährigen Greises hat es noch mit ins Grab genommen. Auf dem kunstvollen Kreuze aber, das Webers Grabstätte schmückt, sind seine Worte eingemeißelt:

\begin{verse}
„Und schlaf ich längst schon unter Friedhofslinden:\\
Das sollst du stets bewahren im Gedächtniß,\\
Als meiner Liebe theuerstes Vermächtniß:\\
Es ist kein Heil als nur im Kreuz zu finden!“
\end{verse}
